\documentclass[11pt]{article}
\usepackage[margin=0.85in]{geometry}
\usepackage{amsmath,amssymb,booktabs,longtable,xurl,hyperref}
\usepackage[T1]{fontenc}
\hypersetup{colorlinks=true,urlcolor=blue,linkcolor=blue}
\setlength{\emergencystretch}{2em}
\title{Restoring CLF Optimization After Zero PCBF Slack}
\author{Pan Dynamics Collision Avoidance Research}
\date{September 30, 2026}
\begin{document}
\maketitle
\section{Confirmed defect and corrected stopping contract}
The preceding implementation at
\path{e2813a5f0d6be6ba7f8a25de4c3850bf28e9a7a6}
returned as soon as a nonlinear trajectory attained zero PCBF safety slack and
satisfied the original hard constraints. This happened both at initialization
and after a primary optimization result. A safe shifted trajectory therefore
bypassed CLF optimization even when its CLF slack was positive. Existing tests
incorrectly required this behavior. The user clarified that completing safety
feasibility must not suppress the second CLF stage.

The corrected single-controller contract is:
\begin{enumerate}
\item Minimize the nonnegative PCBF collision-slack sum until a trajectory has
zero slack and satisfies every original hard constraint.
\item Preserve that zero safety level while optimizing CLF slack, with input
proximity measured from the current linearization trajectory.
\item Stop after a nonlinearly admitted secondary update, or when the current
CLF penalty is already within tolerance of its nonnegative lower bound.
\end{enumerate}
A safe initialization proves that the primary minimum is zero, so the primary
conic solve can be omitted. It does not prove that the CLF stage is complete.
A soft deadline can leave the admitted incumbent as the issued plan, but the
metadata explicitly records incomplete CLF work. No additional controller,
backup selector, road constraint or terminal mode is introduced.

\section{Secondary formulation and nonlinear admission}
Let $\bar U$ denote the current nonlinear rollout inputs, $d_U=U-\bar U$,
$W(e)=\|Le\|_2^2$, and $c=(1-\alpha)W(e_0)$.
The affine first-successor error is $\widehat e_1=e_1+J_ed$.
Instead of a tangent plane of $W$, the conic program retains its full convex
quadratic in the affine error:
\[
\|L\widehat e_1\|_2^2\le c+\rho,\quad \rho\ge0
\quad\Longleftrightarrow\quad
\left\|\begin{bmatrix}2L\widehat e_1\\\rho+c-1\end{bmatrix}\right\|_2
\le\rho+c+1.
\]
The tangent plane omitted a nonnegative quadratic error term and could
understate the local CLF slack. This cone removes that source of underestimation;
it does not make the nonlinear dynamics affine or remove nonlinear validation.
The secondary objective retains the original horizon tracking cost:
\[
\min\quad w_\rho\rho^2+\frac{1}{H}\sum_{i=0}^{H-1}
\left(\|L\widehat e_{i+1}\|_2^2+0.01
\|\operatorname{diag}(\sqrt{w_\delta},\sqrt{w_\beta})d_{u,i}\|_2^2\right).
\]
The existing positive weights are unchanged ($w_\rho=100$ and input weights
one by default). Input magnitude is penalized relative to $\bar U$, not zero.
Both tracking and the CLF reference the original given path. Stored nonlinear
tracking/CLF cost excludes input proximity because that term depends on the
current anchor. It supplies the tie-break between equal actual CLF slacks.

An intermediate simplification removed horizon tracking. Although its complete
fourteen-case campaign remained collision-free, final lateral error reached
74.024 m. That term was therefore restored. Earlier slow whole-tail-cost trials
also lacked adequate scaling and feasible-anchor backtracking; they did not
establish that the tracking term itself should be removed.

At an admitted anchor all collision slacks and terminal elastic variables
are fixed exactly to zero. The conic CLF slack has upper bound
$\rho_{\rm anchor}$. Actual nonlinear rollout admission checks the original
collision, physical, input, slew and terminal conditions. A secondary result
must not increase actual CLF slack. Raw, corrected and secondary trajectories
remain subject to the same candidate comparison.

The CLF penalty alone is nonnegative. If
$w_\rho\rho_{\rm anchor}^2$ is at most the existing $10^{-7}$ tolerance,
its gap to the zero lower bound is certified and no secondary solve is needed.
This does not establish optimality of the combined tracking objective.
The stopping contract requires no further tracking-cost refinement once CLF
dissipation is attained. Otherwise a secondary solve is attempted.
The output distinguishes \texttt{clfStageAttempted},
\texttt{clfStageCompleted}, and \texttt{clfLowerBound}.
Finite suboptimal conic points can be admitted after nonlinear checking;
reported conic exit flags are preserved. Stage completion means one accepted
local SCvx update, not global optimality or nonlinear stationarity.

\subsection{Scaling and an informative flow anchor}
The first complete braking-lead run with the unscaled repair was collision-free
but had 54 unfinished CLF calls and an 84.333 m final lateral error. This was
not accepted as the final implementation. Large residuals made the squared-slack
epigraph numerically difficult. The solver now uses
$s=\max(1,\rho_{\rm anchor})$, $\widehat\rho=\rho/s$, divides squared CLF
quantities by $s$, and divides the entire secondary objective by the positive constant $\max(1,J_{\rm anchor})$, using the anchor tracking/CLF cost.
This preserves the physical feasible set and objective ordering in exact
arithmetic. Solver tolerances apply in the scaled conic units; nonlinear
admission and the no-solve lower-bound test use original units. All six isolated
scaled secondary probes returned normal optimal exit flags.

The old flow initializer used a kinematic steering preference. At 8 m/s its
0.307802 rad initial limit exceeded the front Fiala saturation threshold
0.228534 rad even at zero longitudinal input. On that branch the front-tire
slip derivative is zero. A local CLF solve can then increase steering to reduce
the cosine-projected saturated force, moving in an undesirable direction.
The initializer now uses the same 0.8 force-fraction preference through the
actual adhesion law $|F_y|/F_{y,\max}=1-(1-q)^3$:
\[
\alpha_{\max}=\arctan\left[
\frac{3\mu_fF_{z,f}\sqrt{1-\beta^2}}{C_f}
\big(1-(1-0.8)^{1/3}\big)\right].
\]
It centers the corresponding steering interval at the current front-axle
zero-slip heading. At zero longitudinal input this gives 0.096275 rad.
Actual amplitude and slew clipping takes precedence. This is only a search-seed
preference; no optimizer constraint or vehicle model is changed. In an isolated CLF-only
8 m/s braking-lead first-call comparison, admitted CLF slack decreased from approximately
1.8323 with the saturated seed to 0.0571 with the informative seed.

\subsection{Feasible-anchor backtracking}
The previous trial check tried only full, half and quarter steps before
rebuilding the conic program, and repeatedly attempted endpoint shooting even
from an admitted anchor. Restoring horizon tracking exposed repeated five-second
calls under that rule. The final solve instead halves one direction geometrically
within the existing budget, up to twenty halvings. Every trial still requires
nonlinear safety and nonincreasing actual CLF slack. Endpoint shooting remains
available for infeasible restoration; it is not repeated along a feasible-anchor
line search. A 160-hold 15 m/s crossing probe reduced its final lateral error
from 74.024 m in the CLF-only variant to 18.987 m, with a maximum running call of
0.876963 s. This remains imperfect path recovery, not a convergence proof.

\section{Unchanged-fixture validation}
14 of fourteen eight-second scenarios pass the independent sampled
collision-avoidance audit. Each uses 160 holds of 50 ms at reference speeds
8 or 15 m/s; the given-path constant-speed baseline collides in every fixture.
The configuration, target, baseline and random-seed records match the previous
campaign exactly. Road constraints remain absent and the nominal prediction
margin remains 6 mm. Prefix lengths are 8/16, maximum horizon 512, soft budget
5 s and outer iteration limit 24. No random draws are used.

Measured ODE45 successors feed back into the controller, using relative/absolute
tolerances $10^{-11}/10^{-12}$ and 31 rectangle-audit samples per hold.
Python independently recomputes target motion, baseline penetration and body
separation. All issued plans have exactly zero evaluated safety slack and
original hard residual. The minimum independently replayed clearance is
\textbf{0.048648 m}, above the 6 mm nominal prediction margin. The intermediate CLF-only
replay reached 0.000436 m, below that prediction margin: RK4 prediction and
independent ODE replay are distinct trajectories. Neither
these finite audit samples nor zero nominal slack establishes continuous-time
or uncertain-plant safety.

\begin{center}\small
\begin{tabular}{rlrrrr}
\toprule
Speed & Scenario & Holds & Min gap (m) & Max running (s) & Final $e_y$ (m)\\
\midrule
8 & headOn & 160 & 1.503071 & 0.717269 & -9.797 \\
8 & acceleratingHeadOn & 160 & 1.206900 & 0.775880 & -9.882 \\
8 & brakingLead & 160 & 0.844861 & 2.057584 & -0.868 \\
8 & crossing & 160 & 0.784016 & 0.718353 & -11.688 \\
8 & turningCrossing & 160 & 0.757383 & 1.938550 & -10.948 \\
8 & curvedHeadOn & 160 & 1.662495 & 1.073695 & -12.177 \\
8 & curvedCrossing & 160 & 0.449089 & 2.462262 & -13.833 \\
15 & headOn & 160 & 0.179999 & 0.723179 & -0.998 \\
15 & acceleratingHeadOn & 160 & 0.048648 & 0.982425 & -0.502 \\
15 & brakingLead & 160 & 0.120129 & 2.452237 & -2.946 \\
15 & crossing & 160 & 1.359353 & 0.745458 & -18.987 \\
15 & turningCrossing & 160 & 1.362759 & 1.138676 & -24.482 \\
15 & curvedHeadOn & 160 & 1.744633 & 0.749704 & -24.740 \\
15 & curvedCrossing & 160 & 1.011963 & 1.297733 & -23.581 \\
\bottomrule
\end{tabular}
\end{center}

There are 1616 completed secondary-update calls and 624 lower-bound
terminations. 0 issued calls report incomplete CLF work.
Secondary conic exit-flag counts are \texttt{1: 883, -7: 743}; these are not silently
reclassified as optimal. Actual nonlinear CLF comparisons are checked in the
compact evidence builder. A positive soft CLF slack can remain necessary or
locally unresolved during avoidance. The maximum final lateral error is
24.739862 m, compared with 30.641000 m in the prior safety-only run; this is an
eight-second observation, not a proof of asymptotic path convergence.

\section{Runtime and regression checks}
Across 2226 running calls, the mean is 470.683069 ms, median 113.279000 ms and
maximum \textbf{2.462262 s}. The maximum occurs in 8 m/s curvedCrossing
at $t=2.050000$ s with 213 horizon holds. Initialization is excluded;
its separate maximum is 1.975543 s. There are 1639 running calls exceeding
the 50 ms sample period. The preceding recorded safety-only maximum was
0.216044 s. That value is not representative of the corrected two-stage
algorithm because the earlier code omitted positive-CLF work.

These are actual unprofiled single-thread controller-call wall times, not
solver-reported time. Other MATLAB research processes were active on the host.
The previous campaign is reused, not a fresh isolated-machine comparison.
A 50 ms deadline is not established, and the five-second limit remains soft.

All 290 behavior tests across ten controller/model/geometry/configuration
classes pass in 40.870466 s. Updated cases cover positive-CLF safe seeds,
positive-CLF exact shifted solutions, primary feasibility followed by the
secondary stage, zero-CLF lower-bound stopping, and a deadline that must not
claim CLF completion. MATLAB Code Analyzer finds fourteen sparse-indexing
advisories in the solver and no findings in the other five checked files.
The nine-source-file controller budget is unchanged.

Intermediate 40-hold head-on/crossing smoke tests completed before the final campaign.
An earlier whole-tail-cost prototype was interrupted after repeated deadline
exhaustion and is preserved only as exploratory diagnostic evidence. It is not
part of the final success count. After improving the seed, the previous 24 m primary-restoration unit fixture no longer needed restoration; its 22 m variant preserves that test coverage. Two closer exploratory variants at 20 and 18 m failed to find a continuation within 20 s and remain outside the unchanged fourteen-case campaign. No physical-impossibility conclusion is drawn. The experiment's transverse error diagnostic
was also corrected to use the separate nominal reference instead of the free
terminal reference; lateral position is unaffected, while curved trim-related
velocity/yaw components now reference their intended nominal values.

\section{Reproduction and evidence}
Implementation and derivation are in
\path{controller/solvePredictiveControl.m} and
\path{controller/PCBF_CLF_ARCHITECTURE.md}. The compact bundle
\path{report/RESTORE_CLF_STAGE_20260930/} contains per-call metrics, final, prior safety-only and intermediate CLF-only audits,
fixture records, test results, analyzer findings, and source/raw manifests.
\path{REPRODUCTION.txt} lists exact drivers and configuration.
Full pose streams and prototype logs remain under
\path{/home/zai/.cache/collisionAvoidance/restore-clf-stage-20260930/}.
No solver/vendor tree, unrelated observer work, generated binary or copied
source snapshot is included in the project commit.
\end{document}
